\documentclass[12pt, a4paper]{article}

% ===== Encoding & Sprache =====
\usepackage[utf8]{inputenc}
\usepackage[T1]{fontenc}
\usepackage[ngerman]{babel}

% ===== Typografie =====
\usepackage[scaled]{helvet}
\renewcommand{\familydefault}{\sfdefault}
\usepackage{microtype}
\usepackage{setspace}
\onehalfspacing

% ===== Seitenränder =====
\usepackage[left=3cm, right=3cm, top=3cm, bottom=3cm]{geometry}

% ===== Literaturverzeichnis (Autor-Jahr-Stil) =====
\usepackage[
  backend=biber,
  style=authoryear,
  sorting=nyt,
  dashed=false,
  maxcitenames=3,
  maxbibnames=99,
  uniquelist=false,
  uniquename=false,
  ibidtracker=false,
  pagetracker=false,
  natbib=true
]{biblatex}
\addbibresource{bibliography.bib}

% ===== Absatzformat =====
\setlength{\parindent}{0pt}
\setlength{\parskip}{6pt}

% ===== Hyperlinks =====
\usepackage[hidelinks]{hyperref}

% ===== Mathematik =====
\usepackage{amsmath}
\usepackage{amssymb}
\DeclareMathOperator*{\argmin}{arg\,min}

% ===== Grafiken & Tabellen =====
\usepackage{graphicx}
\usepackage{booktabs}
\usepackage{caption}

% ===== Algorithmen (für die GA-Beschreibung) =====
\usepackage{algorithm}
\usepackage{algpseudocode}
\floatname{algorithm}{Algorithmus}

% ===== Kopf-/Fußzeile =====
\usepackage{fancyhdr}
\pagestyle{fancy}
\fancyhf{}
\rhead{\thepage}
\renewcommand{\headrulewidth}{0pt}

% ===== Titeldaten (hier anpassen) =====
\newcommand{\ArbeitTitel}{Genetic Algorithm für\\[0.4cm]Multi-Objective Trip Planning}
\newcommand{\ArbeitAutorEins}{Marcel Rommel ()}
\newcommand{\ArbeitAutorZwei}{Jessica Noe (8595149)}
\newcommand{\ArbeitModulEins}{Methoden der Künstlichen Intelligenz}
\newcommand{\ArbeitModulZwei}{und Computational Intelligence}
\newcommand{\ArbeitModulNr}{T3M40610}
\newcommand{\ArbeitBetreuer}{Prof. Dr. Heinrich Braun}
\newcommand{\ArbeitBearbeitungszeitraum}{15.10.2026}

\hypersetup{pdftitle={\ArbeitTitel}, pdfauthor={\ArbeitAutorEins, \ArbeitAutorZwei}}

% ---------------------------------------------------------

\begin{document}

% ===== TITELBLATT =====
\begin{titlepage}
  \raggedleft
  \includegraphics[width=8cm]{img/dhbw-cas-logo}\par
  \centering
  \vspace*{1.5cm}

  {\large Seminararbeit}\\[1cm]

  \begin{singlespacing}
  {\LARGE\bfseries \ArbeitTitel}
  \end{singlespacing}
  \vspace{1.5cm}

  {\large am Center for Advanced Studies}\\[0.2cm]
  {\large der Dualen Hochschule Baden-Württemberg}\\[2cm]

  \begin{tabular}{ll}
  	Modul: & \ArbeitModulEins \\
  				& \ArbeitModulZwei \\
  				& \ArbeitModulNr \\
    Verfasser: & \ArbeitAutorEins \\
               & \ArbeitAutorZwei \\
    Abgabedatum: & \ArbeitBearbeitungszeitraum \\
    Dozent: & \ArbeitBetreuer \\
  \end{tabular}

  \vfill

\end{titlepage}

% ===== EHRENWÖRTLICHE ERKLÄRUNG =====
\newpage
\thispagestyle{empty}

{\large\bfseries Ehrenwörtliche Erklärung}\\[1em]

Wir erklären hiermit, dass wir die vorliegende Seminararbeit selbstständig
verfasst und keine anderen als die angegebenen Hilfsmittel benutzt haben.
Alle Stellen, die wörtlich oder sinngemäß aus veröffentlichten oder nicht
veröffentlichten Quellen entnommen wurden, sind als solche kenntlich
gemacht. Die Arbeit wurde in gleicher oder ähnlicher Form noch keiner
anderen Prüfungsbehörde vorgelegt.

\vspace{3em}

\begin{tabular}{@{}p{7cm}@{}}
  \dotfill \\
  Ort, Datum, Unterschrift \ArbeitAutorEins
\end{tabular}

\vspace{3em}

\begin{tabular}{@{}p{7cm}@{}}
  \dotfill \\
  Ort, Datum, Unterschrift \ArbeitAutorZwei
\end{tabular}

\vspace{2em}

{\large\bfseries Erklärung zur Nutzung generativer KI}\\[1em]

Bei der Erstellung dieser Arbeit wurden generative KI-Systeme eingesetzt, namentlich
Claude (Anthropic). Die Nutzung erfolgte in folgendem Umfang: Strukturierung der
Kapitel~\ref{sec:modellierung} und~\ref{sec:loesungsverfahren}, Unterstützung bei der
mathematischen Formalisierung und Notation, Überarbeitung des
LaTeX-Quelltextes, Quellensuche sowie sprachliche Überarbeitung.

Die inhaltlichen Entscheidungen der Modellierung und des Lösungsverfahrens wurden von
den Verfassern getroffen. Sämtliche mithilfe von KI erzeugten Inhalte wurden geprüft,
überarbeitet und inhaltlich verantwortet. Die Verantwortung für die Richtigkeit und
Vollständigkeit der Arbeit einschließlich aller Quellenangaben liegt bei den
Verfassern.

% ===== VERZEICHNISSE =====
\newpage
\pagenumbering{Roman}

\tableofcontents

\newpage
\listoftables
\addcontentsline{toc}{section}{Tabellenverzeichnis}

% ===== HAUPTTEIL =====
% Seitenzahl der Verzeichnisse merken, damit das Literaturverzeichnis
% die roemische Zaehlung fortsetzen kann.
\newcounter{romanpages}
\setcounter{romanpages}{\value{page}}

\newpage
\pagenumbering{arabic}

\section{Einleitung}

\subsection{Problemstellung}

Moderne Fernreisen umfassen häufig mehrere Ziele: Laut einer Studie der European Travel Commission besuchten 92\,\% der befragten Fernreisenden aus den USA, China, Japan und Indien bei ihrer letzten Europareise mehrere Städte, 60\,\% sogar mehr als ein Land. Multi-Destination-Trips sind damit die vorherrschende Reiseform und nicht die Ausnahme \parencite[Vgl.][S.~7, 11]{etcTrackingMultiDestinationTravel2018}.

Mit jedem zusätzlichen Ziel wächst jedoch der Planungsaufwand, denn Reihenfolge der Städte, Aufenthaltsdauern, Transportmittel und Unterkünfte müssen gemeinsam festgelegt werden und beeinflussen sich gegenseitig. Bereits eine um einen Tag veränderte Aufenthaltsdauer verschiebt alle folgenden Etappen und damit auch deren Preise, die unter anderem vom Wochentag abhängen, sowohl bei Flügen \parencite[Vgl.][]{mantinKooWeekendEffect2010} als auch bei Hotelübernachtungen \parencite[Vgl.][]{schamelWeekendMidweekHotel2012}.

Neben Kosten und Komfort rückt dabei zunehmend eine dritte Zielgröße in den Fokus: die Klimawirkung der Reise. Reisebedingte Emissionen stiegen von 982 Millionen Tonnen CO\textsubscript{2} (2005, rund 5\,\% aller anthropogenen Emissionen) auf 1.597 Millionen Tonnen (2016) und sollen bis 2030 auf etwa 1.998 Millionen Tonnen weiter ansteigen \parencite[Vgl.][S.~49]{unwtoTransportrelatedCO2Emissions2019}. Entscheidend ist hier vor allem die Wahl des Transportmittels: Nur rund ein Drittel aller Reisen wurde 2016 mit dem Flugzeug zurückgelegt, auf sie entfiel jedoch die Hälfte der reisebedingten CO\textsubscript{2}-Emissionen \parencite[Vgl.][S.~46]{unwtoTransportrelatedCO2Emissions2019}.

Eine gute Reiseplanung muss diese teilweise konfligierenden Ziele, Kosten, Komfort und CO\textsubscript{2}, folglich gleichzeitig berücksichtigen. Die Zahl möglicher Kombinationen aus Reihenfolge, Aufenthaltsdauern, Verbindungen und Unterkünften wächst jedoch so schnell, dass eine manuelle Bewertung aller Alternativen ausgeschlossen ist.

\subsection{Lösungsansatz}

Das beschriebene Problem vereint zwei klassische Optimierungsprobleme: Die Reihenfolge der Reiseziele entspricht einem Traveling-Salesman-Problem, die Verteilung der begrenzten Ressourcen Reisezeit, Geld und CO\textsubscript{2}-Kontingent auf die einzelnen Stopps einem Rucksackproblem. Beide sind in ihrer Entscheidungsvariante NP-vollständig und damit in der hier betrachteten Optimierungsvariante NP-schwer \parencite[Vgl.][S.~211, 247]{gareyJohnsonComputersIntractability1979}, und der kombinierte Suchraum aus Reihenfolge, Aufenthaltsdauern, Verbindungen und Hotels wächst bereits bei wenigen Zielen so stark, dass eine vollständige Enumeration ausscheidet.

Für solche Problemklassen bleibt nur, auf exakte Lösungen zu verzichten und mit Approximationsverfahren und Metaheuristiken zu arbeiten \parencite[Vgl.][S.~11--12]{eibenSmithIntroductionEvolutionary2015}. Evolutionäre Verfahren eignen sich dafür besonders, da sie in vertretbarer Rechenzeit gute Näherungslösungen liefern und mit gemischten Repräsentationen umgehen können, die sich für Probleme mit unterschiedlich strukturierten Entscheidungen häufig besser eignen als eine einheitliche Kodierung \parencite[Vgl.][S.~49]{eibenSmithIntroductionEvolutionary2015}. In dieser Arbeit wird daher ein genetischer Algorithmus entworfen: Ein Reiseplan wird als Genstring aus Tripeln (Reiseziel, Aufenthaltsdauer und Verbindung) kodiert und über eine gewichtete Fitnessfunktion aus Kosten, Diskomfort und CO\textsubscript{2} bewertet, deren Gewichte die Präferenzen der reisenden Person abbilden. Die Wahl der Unterkunft wird nicht mitevolviert, sondern bei der Dekodierung eines Genstrings ergänzt. Problemspezifische Mutations- und Crossover-Operatoren berücksichtigen dabei die unterschiedliche Struktur der Genabschnitte, insbesondere den Permutationscharakter der Reihenfolge.

Der Schwerpunkt dieser Arbeit liegt auf der Modellierung des Problems. Die prototypische Umsetzung veranschaulicht ausgewählte Komponenten des Verfahrens.

\section{Modellierung}
\label{sec:modellierung}

\subsection{Formalisierung des Planungsproblems}
\label{subsec:formalisierung}

Die Planung einer Mehrziel-Reise wird im Folgenden als kombinatorisches Optimierungsproblem aufgefasst. Gegeben ist eine Menge $Z = \{1, \dots, n\}$ von Reisezielen, die ausgehend von einem festen Startort $z_0$ jeweils genau einmal besucht werden sollen, bevor die Reise wieder an $z_0$ endet. Der Reisezeitraum umfasst $T$ Tage, innerhalb derer sich Aufenthalte und Ortswechsel verteilen.

Tabelle~\ref{tab:eingangsgroessen} fasst die Größen zusammen, die als gegeben vorausgesetzt werden.

\begin{table}[htbp]
  \centering
  \caption{Eingangsgrößen des Planungsproblems}
  \label{tab:eingangsgroessen}
  \begin{tabular}{@{}lp{10.5cm}@{}}
    \toprule
    Symbol & Bedeutung \\
    \midrule
    $Z = \{1,\dots,n\}$ & Menge der zu besuchenden Reiseziele \\
    $z_0$ & Startort, an dem die Reise beginnt und endet \\
    $T$ & Länge des Reisezeitraums in Tagen \\
    $d_{\min}(z)$ & Mindestaufenthaltsdauer am Ziel $z$ \\
    $V(a,b)$ & Verfügbare Verkehrsverbindungen für den Übergang von $a$ nach $b$ \\
    $H(z)$ & Vorausgewählte Hotels am Ziel $z$ \\
    $E$ & Distanzmatrix mit den Entfernungen zwischen den Orten \\
    $B_{\max}$ & Budgetobergrenze für die gesamte Reise \\
    $C_{\max}$ & CO\textsubscript{2}-Kontingent für die gesamte Reise \\
    \bottomrule
  \end{tabular}
\end{table}

Drei dieser Größen bedürfen einer Erläuterung. Die Mindestaufenthaltsdauer $d_{\min}(z)$ bildet ab, dass sich ein Ziel unterhalb einer gewissen Verweildauer nicht sinnvoll erschließen lässt. Eine Großstadt mit mehreren zentralen Sehenswürdigkeiten rechtfertigt den Anreiseaufwand erst ab zwei bis drei Übernachtungen, während für einen kleineren Ort eine Nacht genügen kann. Die Mengen $V(a,b)$ enthalten keine abstrakten Relationen zwischen Orten, sondern konkrete Verbindungen mit Verkehrsmittel, Abfahrtszeitfenster, Fahrtdauer, Preis und CO\textsubscript{2}-Ausstoß. Analog umfasst $H(z)$ eine Vorauswahl von Hotels je Ziel mit tagesabhängigen Übernachtungspreisen und einer Komfortbewertung.

Gesucht ist ein vollständiger Reiseplan, der Besuchsreihenfolge, Aufenthaltsdauer je Stopp, Verbindung je Etappe und Hotel je Stopp gemeinsam festlegt und dabei die Zielgrößen Kosten, Komfort und CO\textsubscript{2}-Ausstoß bestmöglich ausbalanciert.

Die Optimierung findet ausschließlich innerhalb vorab bereitgestellter Optionen statt. Sämtliche Preis-, Emissions- und Verbindungsdaten werden im Vorfeld erhoben und als Nachschlagetabellen vorgehalten, denn ein evolutionärer Algorithmus bewertet je Lauf mehrere tausende Kandidaten gegen einen unveränderlichen Datenstand.

Naheliegend wäre, dieses Problem zu zerlegen und nacheinander zu lösen: zuerst die Route, dann die Aufenthaltsdauern, zuletzt Transport und Unterkunft. Eine solche sequenzielle Behandlung scheitert jedoch an der wechselseitigen Abhängigkeit der Teilentscheidungen. Ein zusätzlicher Übernachtungstag am ersten Stopp verlegt jede weitere Fahrt und jede weitere Hotelnacht um einen Tag und verändert damit deren tagesabhängige Preise. Eine Route, die unter festen Aufenthaltsdauern optimal erscheint, kann bei veränderter Zeitaufteilung deutlich schlechter abschneiden. Die Entscheidungen müssen daher gemeinsam optimiert werden.

\subsection{Entscheidungsvariablen und Repräsentation}
\label{subsec:repraesentation}

Die Repräsentation einer Lösung als Genstring legt fest, welche Entscheidungen der Algorithmus überhaupt variieren kann, und bestimmt damit zugleich, welche Mutations- und Rekombinationsoperatoren sinnvoll definierbar sind. Sie gilt zugleich als einer der schwierigsten Teile beim Entwurf eines evolutionären Algorithmus \parencite[Vgl.][S.~49]{eibenSmithIntroductionEvolutionary2015} und ist damit die zentrale Modellierungsentscheidung dieser Arbeit.

Ein Reiseplan wird als Tripel $G = (\pi, d, t)$ dargestellt, dessen Komponenten die drei frei zu wählenden Entscheidungsdimensionen abbilden:

Die \textbf{Besuchsreihenfolge} $\pi = (\pi_1, \dots, \pi_n)$ ist eine Permutation der Zielindizes $\{1, \dots, n\}$. $\pi_i$ bezeichnet dasjenige Reiseziel, das an Position $i$ der Route besucht wird. Die Permutationseigenschaft stellt sicher, dass jedes Ziel genau einmal vorkommt, und schließt damit doppelte Besuche und Auslassungen bereits auf Repräsentationsebene aus.

Die \textbf{Aufenthaltsdauer} $d = (d_1, \dots, d_n)$ gibt mit $d_i \in \mathbb{N}$ die Anzahl der Übernachtungen am Stopp $\pi_i$ an. Die Wahl ganzer Zahlen bildet ab, dass Hotelübernachtungen und Fahrpläne nur tageweise disponierbar sind. Eine kontinuierliche Modellierung hätte hier keine Entsprechung in der Anwendung.

Die \textbf{Transportwahl} $t = (t_1, \dots, t_{n+1})$ umfasst $n+1$ Einträge, da neben den $n$ Anreisen zu den Stopps auch die Rückreise von $\pi_n$ zum Startort $z_0$ zu treffen ist. Dabei bezeichnet $t_i$ für $i \leq n$ die Verbindung für den Übergang von $\pi_{i-1}$ nach $\pi_i$ (mit $\pi_0 = z_0$) und $t_{n+1}$ die Rückreise.

Für die Transportwahl wird nicht eine konkrete Verbindung kodiert, sondern eine Kombination aus Verkehrsmittel und Abfahrtszeitfenster:
\begin{equation}
t_i \in M \times S, \qquad M = \{\text{Bahn}, \text{Bus}, \text{Flug}\}, \quad S = \{\text{Vormittag}, \text{Nachmittag}, \text{Nacht}\}
\end{equation}

Diese Diskretisierung folgt der Struktur der Entscheidung, wie sie eine reisende Person tatsächlich trifft. Zwischen zwei Zügen um 9:05 und um 9:40 Uhr besteht aus Planungssicht kein wesentlicher Unterschied, zwischen einem Vormittagszug und einem Nachtbus dagegen ein erheblicher. Ein Genstring, der jede einzelne Abfahrtszeit als eigenen Wert führte, würde dem Algorithmus vorwiegend bedeutungslose Variation anbieten. Die Zeitfenster-Kodierung sorgt dafür, dass benachbarte Genwerte auch anwendungsseitig benachbarte Lösungen bezeichnen, was Voraussetzung für die in Abschnitt~\ref{subsec:mutation} beschriebene Nachbarschaftsheuristik ist. Welche konkrete Verbindung innerhalb des gewählten Fensters gebucht wird, entscheidet die Dekodierfunktion (Abschnitt~\ref{subsec:dekodierung}).

Tabelle~\ref{tab:entscheidungsvariablen} fasst die Entscheidungsvariablen zusammen.

\begin{table}[htbp]
  \centering
  \caption{Entscheidungsvariablen des Genstrings}
  \label{tab:entscheidungsvariablen}
  \begin{tabular}{@{}llll@{}}
    \toprule
    Variable & Bedeutung & Wertebereich & Länge \\
    \midrule
    $\pi_i$ & Ziel an Routenposition $i$ & Permutation über $\{1,\dots,n\}$ & $n$ \\
    $d_i$   & Übernachtungen am Stopp $\pi_i$ & $d_i \in \mathbb{N}$, $d_i \geq d_{\min}(\pi_i)$ & $n$ \\
    $t_i$   & Verkehrsmittel und Zeitfenster & $M \times S$, $|M \times S| = 9$ & $n+1$ \\
    \bottomrule
  \end{tabular}
\end{table}

Der Genstring ist damit heterogen aufgebaut und vereint einen Permutationsabschnitt, einen ganzzahligen Abschnitt und einen kategorialen Abschnitt. Solche gemischten Repräsentationen gelten als die natürlichere Wahl, sobald sich die Bestandteile eines Problems nicht sinnvoll in eine einheitliche Form bringen lassen \parencite[Vgl.][S.~49]{eibenSmithIntroductionEvolutionary2015}.

Als durchgängiges Beispiel dient eine Reise mit $n = 4$ Zielen über $T = 14$ Tage. Die Ziele seien Paris (1), Berlin (2), Rom (3) und Barcelona (4). Der Genstring
\begin{align*}
G = \big(&(2,\, 1,\, 3,\, 4),\; (2,\, 3,\, 4,\, 5),\\
         &(\text{Bahn/Vm},\, \text{Bahn/Nm},\, \text{Flug/Vm},\, \text{Flug/Nm},\, \text{Flug/Vm})\big)
\end{align*}
kodiert eine Reise, die am Vormittag mit der Bahn nach Berlin führt, wo zwei Nächte verbracht werden. Es folgen Paris mit drei Nächten, erreicht über eine Nachmittagsverbindung der Bahn, anschließend Rom mit vier Nächten über einen Vormittagsflug und schließlich Barcelona mit fünf Nächten über einen Nachmittagsflug. Die Rückreise zum Startort erfolgt mit einem Vormittagsflug. Die Abkürzungen Vm und Nm stehen dabei für die Zeitfenster Vormittag und Nachmittag. 

\subsection{Dekodierfunktion}
\label{subsec:dekodierung}

Der Genstring ist der Genotyp und enthält ausschließlich Indizes und Kennungen. Der Phänotyp ist der fertige Reiseplan mit Kalendertagen, gebuchten Verbindungen und namentlich benannten Hotels. Die Dekodierfunktion $D$ bildet den einen auf den anderen ab und ergänzt dabei alle Entscheidungen, die bewusst nicht kodiert wurden. Sie arbeitet deterministisch, sodass derselbe Genstring stets denselben Reiseplan und damit denselben Fitnesswert ergibt.

Die Dekodierung verläuft entlang der Route und führt einen Tageszähler mit. Der Ankunftstag $a_i$ am Stopp $\pi_i$ ergibt sich aus dem Abreisetag vom Vorgänger zuzüglich eines Versatzes, der vom gewählten Zeitfenster abhängt:
\begin{equation}
a_i = a_{i-1} + d_{i-1} + \delta(t_i), \qquad
\delta(t_i) = \begin{cases} 0 & \text{Vormittag oder Nachmittag} \\ 1 & \text{Nacht} \end{cases}
\end{equation}

Nachtverbindungen kommen demnach am Folgetag an und ersetzen dabei eine Übernachtung. Das bildet einen realen Zielkonflikt ab, denn ein Nachttransfer spart Hotelkosten und einen Reisetag, geht aber zulasten des Komforts.

Anschließend werden die beiden ausgeklammerten Entscheidungen ergänzt. Für die Unterkunft wählt der Decoder unter den Hotels des Ziels dasjenige mit dem geringsten gewichteten Beitrag aus Übernachtungspreis und Diskomfort:
\begin{equation}
h_i = \argmin_{h \in H(\pi_i)} \Big( w_1 \sum_{\tau = a_i}^{a_i + d_i - 1} p_h(\tau) \;+\; w_2 \, u(h) \, d_i \Big)
\end{equation}

Dabei bezeichnet $p_h(\tau)$ den Übernachtungspreis von Hotel $h$ am Tag $\tau$ und $u(h)$ den Diskomfort pro Nacht. Analog wählt der Decoder innerhalb der durch $t_i$ festgelegten Klasse aus Verkehrsmittel und Zeitfenster die konkrete Verbindung mit dem geringsten gewichteten Beitrag aus Preis, Reisediskomfort und Emissionen.

Diese Auslagerung ist verlustfrei, und zwar aus einem strukturellen Grund. Die Zielfunktion setzt sich, wie Abschnitt~\ref{subsec:zielfunktion} zeigt, additiv aus den Beiträgen der einzelnen Stopps und Etappen zusammen. Die Wahl des Hotels am Stopp $\pi_i$ beeinflusst ausschließlich den Beitrag dieses Stopps. Sie verändert weder den Zeitplan noch die Emissionen noch die an anderen Stopps verfügbaren Optionen. Bei gegebenem Genstring zerfällt die Hotelwahl damit in $n$ voneinander unabhängige Teilprobleme, deren jeweiliges Minimum unmittelbar ablesbar ist. Das stopweise Minimum ist folglich zugleich das globale Minimum, und ein evolutionärer Suchprozess über die Hotelindizes könnte kein besseres Ergebnis finden, sondern nur langsamer zum selben gelangen.

Für die Verbindungen gilt dasselbe Argument innerhalb einer Klasse, da das Zeitfenster den Ankunftstag bereits festlegt. Es gilt jedoch nicht für die Klasse selbst, denn Verkehrsmittel und Zeitfenster bestimmen über $\delta(t_i)$ den weiteren Zeitplan. Genau deshalb bleibt $t$ im Genstring, während $h$ ihn verlässt.

\subsection{Zielfunktion}
\label{subsec:zielfunktion}

Die Zielfunktion bewertet einen dekodierten Reiseplan und macht damit Lösungen vergleichbar. Sie ist der Punkt, an dem die Planungsqualität objektiviert wird, und zugleich die Grundlage jeder Verbesserung durch Mutation oder Rekombination. Sie wird im Folgenden als Zielfunktion $F(G)$ geführt und minimiert, ein kleinerer Wert entspricht also einem besseren Reiseplan.

Die drei Teilziele werden zunächst getrennt berechnet. Sei $v_i$ die vom Decoder gewählte Verbindung der Etappe $i$ und $h_i$ das Hotel am Stopp $\pi_i$. Mit $c(v)$, $e(v)$ und $u(v)$ für Preis, Emissionen und Diskomfort einer Verbindung gilt:
\begin{align}
K(G) &= \sum_{i=1}^{n+1} c(v_i) \;+\; \sum_{i=1}^{n} \sum_{\tau = a_i}^{a_i + d_i - 1} p_{h_i}(\tau) \\
C(G) &= \sum_{i=1}^{n+1} e(v_i) \\
D(G) &= \sum_{i=1}^{n+1} u(v_i) \;+\; \sum_{i=1}^{n} u(h_i)\, d_i
\end{align}

Der Diskomfort einer Verbindung fasst mehrere Eigenschaften zusammen, die eine Reise als anstrengend erscheinen lassen:
\begin{equation}
u(v) = \lambda_1 \cdot \text{Fahrtdauer}(v) + \lambda_2 \cdot \text{Umstiege}(v) + \lambda_3 \cdot \mathbf{1}[\,v \text{ ist Nachtverbindung}\,]
\end{equation}

Die Zusammenführung der drei Teilziele erfolgt als gewichtete Summe. Dabei stellt sich ein Problem, das eine reine Gewichtung nicht löst: Die Teilziele werden in unterschiedlichen Einheiten gemessen, nämlich in Euro, in Kilogramm und in einer dimensionslosen Diskomfortpunktzahl. Eine Gewichtung von $w_1 = w_2 = w_3 = \tfrac{1}{3}$ wäre in dieser Form bedeutungslos, da sie faktisch von den Größenordnungen der Einheiten bestimmt würde. Die Teilziele werden deshalb zuvor normiert. Als Referenz dient der Idealpunkt, also der Vektor aus den Bestwerten, die jedes Teilziel für sich allein erreichen kann. Diese Bestwerte werden in der Vorverarbeitung greedy angenähert, indem je eine Reise erzeugt wird, die ausschließlich das betreffende Einzelziel optimiert:
\begin{equation}
F(G) = w_1 \frac{K(G)}{K_{\text{ref}}} + w_2 \frac{D(G)}{D_{\text{ref}}} + w_3 \frac{C(G)}{C_{\text{ref}}},
\qquad w_1 + w_2 + w_3 = 1,\; w_j \geq 0
\end{equation}

Die gewichtete Summe überführt das Mehrzielproblem in ein Einzielproblem. Dieses Vorgehen wird als Skalarisierung bezeichnet. Der Preis dafür ist, dass je Gewichtung nur eine einzelne Lösung entsteht und dass die Gewichte bereits festgelegt werden müssen, bevor bekannt ist, welche Lösungen überhaupt erreichbar sind \parencite[Vgl.][S.~196]{eibenSmithIntroductionEvolutionary2015}. Für die vorliegende Anwendung ist das vertretbar, denn die reisende Person benötigt am Ende genau einen Reiseplan und kann ihre Präferenz über die Gewichte ausdrücken. Alternativ ließe sich eine Paretofront bestimmen und der Person zur Auswahl vorlegen, was den Rechenaufwand erhöht und die Entscheidung verlagert statt sie zu treffen.

Bemerkenswert an der obigen Form ist die additive Struktur. $F$ zerfällt in Beiträge einzelner Etappen und Stopps. Für die Bewertung mutierter Individuen bedeutet sie zudem, dass nicht immer der gesamte Plan neu berechnet werden muss. Eine Mutation auf $t_i$ innerhalb desselben Zeitfensters verändert nur den Beitrag einer einzigen Etappe. Eine Mutation auf $d_i$ verschiebt dagegen alle nachfolgenden Kalendertage und erzwingt damit die Neuberechnung des restlichen Plans.

\subsection{Nebenbedingungen}
\label{subsec:nebenbedingungen}

Nebenbedingungen grenzen die zulässigen Reisepläne von den unzulässigen ab. Ein Teil von ihnen muss gar nicht eigens geprüft werden, weil die Repräsentation sie bereits ausschließt. Dass jedes Ziel genau einmal besucht wird, folgt unmittelbar aus der Permutationseigenschaft von $\pi$, und dass jeder Stopp ein Hotel erhält, aus der Konstruktion der Dekodierfunktion. Nebenbedingungen auf Repräsentationsebene aufzulösen ist der wirksamste Umgang mit ihnen, da unzulässige Lösungen so überhaupt nicht entstehen können.

Als harte Bedingungen verbleiben die Mindestaufenthaltsdauer und die Zulässigkeit von Flügen. Für Letztere genügt die reine Distanzgrenze nicht, da sie kurze Strecken mit schlechter Bahnanbindung ausschließt, auf denen ein Flug durchaus sinnvoll ist. Ein Flug ist daher zulässig, wenn die Entfernung eine Schwelle überschreitet oder die beste bodengebundene Alternative erheblich länger dauert:
\begin{equation}
E(\pi_{i-1}, \pi_i) > 500\,\text{km} \quad\lor\quad \big( \vartheta_{\text{Boden}}(\pi_{i-1}, \pi_i) - \vartheta_{\text{Flug}}(\pi_{i-1}, \pi_i) \big) > \Delta\vartheta
\end{equation}

Dabei bezeichnet $\vartheta$ die jeweils kürzeste Reisezeit und $\Delta\vartheta$ einen Schwellwert von beispielsweise vier Stunden. Diese Bedingung wird nicht während der Optimierung geprüft, sondern beim Aufbau der Verbindungsmengen: Für Etappen, die sie verletzen, enthält $V(a,b)$ schlicht keine Flugverbindungen. Auch hier entsteht die Unzulässigkeit damit erst gar nicht. Kodiert ein Genstring dennoch eine leere Klasse aus Verkehrsmittel und Zeitfenster, etwa nach einer Mutation von $\pi$, so wird $t_i$ auf die nächstgelegene belegte Klasse zurückgeführt. Die Mindestaufenthaltsdauer wird analog behandelt, indem die Operatoren $d_i$ nach unten auf $d_{\min}(\pi_i)$ begrenzen.

Anders verhält es sich mit den drei globalen Obergrenzen für Reisedauer, Budget und Emissionen. Sie werden als weiche Bedingungen über Straftermen abgebildet, die die Zielfunktion ergänzen:
\begin{equation}
P(G) = \rho_T \max\!\big(0, \Theta(G) - T\big) + \rho_B \frac{\max\!\big(0, K(G) - B_{\max}\big)}{K_{\text{ref}}} + \rho_C \frac{\max\!\big(0, C(G) - C_{\max}\big)}{C_{\text{ref}}}
\end{equation}

Hier ist $\Theta(G)$ die Gesamtdauer der Reise einschließlich der Reisetage. Bewertet wird anschließend $F(G) + P(G)$. Die Strafterme sind auf dieselben Referenzwerte normiert wie die Teilziele, damit sie mit ihnen vergleichbar bleiben.

Die Entscheidung gegen harte Grenzen ist bewusst getroffen. Würden Budget und CO\textsubscript{2}-Kontingent hart erzwungen, zerfiele der Suchraum in zulässige und unzulässige Bereiche, und der Algorithmus könnte nicht mehr durch eine leicht unzulässige Lösung hindurch zu einer besseren zulässigen gelangen. Ein Plan, der das Budget um wenige Euro überschreitet, wäre zudem schlagartig wertlos, obwohl er einem zulässigen Plan sehr ähnlich ist. Die Straftermen erhalten stattdessen die Richtungsinformation: Je weiter eine Lösung über der Grenze liegt, desto schlechter wird sie bewertet. Die Faktoren $\rho$ müssen dabei so groß gewählt werden, dass sich eine Überschreitung im Optimum nicht auszahlt, also größer als der maximal erzielbare Gewinn je Einheit Überschreitung. Fallen sie zu klein aus, liegt das Optimum der erweiterten Zielfunktion außerhalb des zulässigen Bereichs, fallen sie zu groß aus, nähert sich das Verhalten wieder dem einer harten Grenze an.

\section{Lösungsverfahren}
\label{sec:loesungsverfahren}

\subsection{Ablauf des genetischen Algorithmus}
\label{subsec:ablauf}

Das Verfahren folgt der generationsbasierten Grundform eines genetischen Algorithmus. Eine Population von Genstrings wird bewertet, aus ihr werden Eltern ausgewählt, und aus diesen entstehen durch Rekombination und Mutation Nachkommen, die die nächste Generation bilden. Übernommen werden ausschließlich die Nachkommen sowie eine kleine Zahl der besten Individuen der Vorgängergeneration. Alle übrigen Eltern entfallen, wodurch die Population stets dieselbe Größe behält.

Algorithmus~\ref{alg:ga} zeigt den Ablauf. Jede Bewertung schließt die Dekodierung nach Abschnitt~\ref{subsec:dekodierung} ein, denn ein Genstring lässt sich erst als vollständiger Reiseplan beurteilen. Bewertet wird dabei stets die um die Straftermen erweiterte Zielfunktion $F(G) + P(G)$.

\begin{algorithm}[htbp]
  \caption{Genetischer Algorithmus zur Reiseplanung}
  \label{alg:ga}
  \begin{algorithmic}[1]
    \Require Eingangsdaten, Gewichte $w_1, w_2, w_3$, Populationsgröße $\mu$, Elitenzahl $\kappa$
    \State erzeuge Startpopulation $P$ aus $\mu$ Individuen
    \State dekodiere und bewerte alle $G \in P$
    \While{kein Abbruchkriterium erfüllt}
      \State $P' \gets$ die $\kappa$ besten Individuen aus $P$
      \While{$|P'| < \mu$}
        \State wähle Eltern aus $P$ durch rangbasierte Selektion
        \State erzeuge Nachkommen durch Rekombination
        \State mutiere den Nachkommen
        \State stelle die Zulässigkeit der veränderten Gene wieder her
        \State dekodiere und bewerte den Nachkommen
        \State füge den Nachkommen zu $P'$ hinzu
      \EndWhile
      \State $P \gets P'$
    \EndWhile
    \State \Return bestbewertetes gefundenes Individuum
  \end{algorithmic}
\end{algorithm}

Die Startpopulation wird überwiegend zufällig erzeugt, um den Lösungsraum breit abzudecken. Ergänzend lassen sich einzelne Individuen greedy konstruieren, etwa über eine Nearest-Neighbor-Route oder über die jeweils günstigsten Hotels. Solche Startlösungen beschleunigen die frühen Generationen, erhöhen aber zugleich die Gefahr, dass sich die Population vorzeitig auf deren Umgebung zusammenzieht. Da Selektion und Mutation den Verlauf ohnehin dominieren, fällt der Einfluss der Startpopulation auf das Endergebnis gering aus, weshalb hier ein überwiegend zufälliger Start gewählt wird.

Tabelle~\ref{tab:parameter} fasst die Parameter des Verfahrens zusammen. Ihre endgültigen Werte lassen sich nicht analytisch bestimmen, sondern nur experimentell ermitteln, weshalb Bereiche angegeben sind.

\begin{table}[htbp]
  \centering
  \caption{Parameter des Verfahrens}
  \label{tab:parameter}
  \begin{tabular}{@{}llp{6.5cm}@{}}
    \toprule
    Parameter & Bereich & Bemerkung \\
    \midrule
    Populationsgröße $\mu$ & 50 bis 500 & Kompromiss aus Rechenzeit und Vielfalt \\
    Elitenanteil $\kappa/\mu$ & 5 bis 20\,\% & sichert das bisher Beste gegen Verlust \\
    Mutationsrate & 10 bis 30\,\% & je Gen und Nachkomme \\
    Max. Generationen & 200 & feste Obergrenze \\
    Zeitbudget & 600 s & feste Obergrenze \\
    Konvergenzschwelle & 1\,\% über 20 Generationen & adaptives Kriterium \\
    \bottomrule
  \end{tabular}
\end{table}

\subsection{Mutation und Nachbarschaft}
\label{subsec:mutation}

Die Mutation ist der Operator, der lokale Verbesserung überhaupt erst ermöglicht, und ihre Gestaltung entscheidet maßgeblich über die Güte des Verfahrens. Leitgedanke ist dabei der Begriff der Nachbarschaft: Eine Mutation soll einen Reiseplan erzeugen, der dem Elternplan ähnlich ist, sodass der Unterschied in der Zielfunktion auf wenige veränderte Bestandteile zurückgeht. Verändert eine Mutation zu viel auf einmal, verhält sich der Nachkomme wie eine zufällig neue Lösung, und die Bewertung trägt keine verwertbare Information mehr über die Richtung der Verbesserung. Da der Genstring aus drei strukturell verschiedenen Abschnitten besteht, benötigt jeder von ihnen eigene Operatoren. Tabelle~\ref{tab:mutationen} gibt einen Überblick.

\begin{table}[htbp]
  \centering
  \caption{Mutationsoperatoren je Genabschnitt}
  \label{tab:mutationen}
  \begin{tabular}{@{}llp{7cm}@{}}
    \toprule
    Abschnitt & Operator & Wirkung \\
    \midrule
    $\pi$ & Swap & vertauscht zwei Positionen, verändert bis zu vier Übergänge \\
    $\pi$ & Insert & entnimmt ein Ziel und setzt es an anderer Stelle ein, verändert drei Übergänge \\
    $\pi$ & Teilstring-Inversion & kehrt einen zusammenhängenden Abschnitt um, verändert zwei Übergänge \\
    $d$ & Gauß-Mutation & verändert eine Aufenthaltsdauer um einen gerundet normalverteilten Betrag \\
    $d$ & Tagestransfer & verschiebt eine Nacht von einem Stopp zu einem anderen \\
    $t$ & Fensterwechsel & wechselt das Zeitfenster bei gleichem Verkehrsmittel \\
    $t$ & Verkehrsmittelwechsel & wechselt das Verkehrsmittel bei gleichem Zeitfenster \\
    \bottomrule
  \end{tabular}
\end{table}

Die drei Operatoren auf $\pi$ unterscheiden sich darin, wie stark sie die bestehende Routenstruktur auflösen. Die Teilstring-Inversion ist der schonendste von ihnen, da sie die Reihenfolge innerhalb des umgekehrten Abschnitts erhält und lediglich zwei Übergänge neu bildet. Swap greift am stärksten ein, weil beide getauschten Ziele jeweils neue Vorgänger und Nachfolger erhalten. Diese Abstufung erlaubt es, die Eingriffstiefe zu steuern: In frühen Generationen sind stärkere Eingriffe sinnvoll, später überwiegen die schonenden. Zusätzlich lässt sich die Auswahl der Positionen steuern, statt sie gleichverteilt zu ziehen. Bevorzugt man Positionen, deren aktuelle Übergänge besonders lange Distanzen aufweisen, so greift die Mutation dort an, wo die Route am ehesten ungünstig ist. Das entspricht der aus dem Traveling-Salesman-Problem bekannten Nachbarschaftsheuristik und erhöht den Anteil aussichtsreicher Nachkommen.

Auf $d$ wirkt eine gerundete Gauß-Mutation, die kleine Änderungen deutlich wahrscheinlicher macht als große und anschließend auf $d_{\min}(\pi_i)$ begrenzt wird. Die Standardabweichung $\sigma$ wird im Verlauf des Laufs verringert, sodass das Verfahren zunächst grob und später fein sucht. Der Tagestransfer ergänzt diesen Operator um eine Variante, die die Gesamtreisedauer unverändert lässt, indem sie einem Stopp eine Nacht nimmt und einem anderen eine gibt. Eine reine Gauß-Mutation verändert dagegen stets auch die Gesamtdauer und erzeugt dadurch leicht Individuen, die den Strafterm für den Reisezeitraum auslösen.

Auf $t$ folgt die Nachbarschaft aus der in Abschnitt~\ref{subsec:repraesentation} gewählten Kodierung. Da ein Genwert aus Verkehrsmittel und Zeitfenster besteht, ist ein benachbarter Wert einer, der genau eine dieser beiden Eigenschaften ändert. Der Wechsel des Zeitfensters verschiebt die Etappe innerhalb des Tages und wirkt vor allem auf Preis und Komfort, während der Wechsel des Verkehrsmittels zusätzlich die Emissionen deutlich verändert. Ein Sprung auf einen beliebigen anderen Wert bleibt mit geringer Wahrscheinlichkeit möglich, damit auch entfernte Bereiche des Suchraums erreichbar bleiben.

Die Stärke der Mutation lässt sich über die Erfolgsquote steuern. Führt nur ein kleiner Teil der Nachkommen zu einer Verbesserung, ist der Eingriff zu grob gewählt, führt ein sehr großer Teil zu Verbesserungen, ist er zu vorsichtig und der Fortschritt je Generation bleibt gering. Die bekannte Ein-Fünftel-Regel gibt hierfür einen Richtwert, nach dem der Fortschritt am größten ist, wenn etwa ein Fünftel der Mutationen eine Verbesserung erzielt \parencite[Vgl.][S.~133]{eibenSmithIntroductionEvolutionary2015}. Sie lässt sich unmittelbar zur Anpassung von $\sigma$ und der Eingriffstiefe auf $\pi$ verwenden. Nach jeder Mutation wird abschließend die Zulässigkeit wiederhergestellt, indem $d_i$ auf die Mindestdauer angehoben und $t_i$ auf eine belegte Klasse zurückgeführt wird.


\subsection{Rekombination}
\label{subsec:rekombination}

Die Rekombination kann auf dem Permutationsabschnitt nicht in ihrer einfachen Form angewandt werden. Ein Ein-Punkt-Crossover, das die vordere Hälfte des einen und die hintere Hälfte des anderen Elternteils zusammensetzt, erzeugt Routen, in denen Ziele doppelt vorkommen und andere fehlen. Der Nachkomme wäre kein gültiger Reiseplan.

Verwendet wird daher ein Order Crossover. Dabei wird ein zusammenhängender Abschnitt der Route des ersten Elternteils unverändert übernommen, und die verbleibenden Positionen werden mit den noch fehlenden Zielen in der Reihenfolge aufgefüllt, in der sie beim zweiten Elternteil auftreten. Der Nachkomme ist so stets eine gültige Permutation und erbt von beiden Eltern einen Teil der Reihenfolgeinformation.

Für die übrigen Genabschnitte stellt sich die Frage, woran sie gebunden sind. Die Aufenthaltsdauer gehört inhaltlich zum Reiseziel und nicht zur Routenposition, denn wie lange sich ein Aufenthalt lohnt, hängt vom Ort ab. Sie wird deshalb gemeinsam mit dem Ziel vererbt, das sie beschreibt. Die Transportwahl gehört dagegen zum Übergang zwischen zwei Zielen, und dieser Übergang entsteht beim Order Crossover teilweise neu. Für neu gebildete Übergänge wird $t_i$ daher nicht vererbt, sondern auf eine belegte Klasse gesetzt und der weiteren Optimierung durch Mutation überlassen.

Ob die Rekombination in diesem Verfahren überhaupt einen Beitrag leistet, ist offen. Sie ist dann wirksam, wenn die Eltern komplementäre Stärken besitzen, etwa wenn einer eine günstige Route und der andere eine geschickte Zeitaufteilung gefunden hat. Bei stark lokal geprägten Problemen kann eine rein mutationsbasierte Suche jedoch ebenso gute Ergebnisse liefern, da die Rekombination bestehende gute Teilstrukturen auch zerstören kann. Der Entwurf hält die Rekombination deshalb abschaltbar, sodass ihr Beitrag experimentell durch den Vergleich von Läufen mit und ohne sie bestimmt werden kann.

\subsection{Selektion und Abbruch}
\label{subsec:selektion}

Die Auswahl der Eltern erfolgt rangbasiert. Dazu wird die Population nach ihrer Bewertung sortiert, und die Auswahlwahrscheinlichkeit eines Individuums ergibt sich allein aus seinem Rang in dieser Sortierung, nicht aus dem Zahlenwert seiner Bewertung. Der Grund liegt in der Zielfunktion selbst. Durch die Straftermen aus Abschnitt~\ref{subsec:nebenbedingungen} sind die Bewertungen nach oben unbeschränkt, denn ein Individuum mit stark überschrittenem Budget erreicht einen um Größenordnungen schlechteren Wert als alle übrigen. Eine fitnessproportionale Auswahl würde durch solche Ausreißer verzerrt, in frühen Generationen ebenso wie später, wenn sich die Bewertungen der guten Individuen einander annähern und kaum noch unterscheidbar sind. Die rangbasierte Auswahl ist gegenüber beiden Effekten unempfindlich, da sie nur die Ordnung der Individuen verwendet. Der Selektionsdruck lässt sich dabei über das Gefälle der Wahrscheinlichkeiten zwischen bestem und schlechtestem Rang einstellen.

Ergänzt wird die Selektion durch Elitismus. Die besten $\kappa$ Individuen einer Generation werden unverändert in die nächste übernommen, sodass die beste bisher gefundene Lösung nicht durch Mutation oder Rekombination verlorengehen kann. Der Wert der besten Lösung verbessert sich dadurch über die Generationen hinweg monoton.

Die Optimierung endet, sobald eines von drei Kriterien eintritt. Zwei davon begrenzen den Ressourcenverbrauch fest, nämlich eine Obergrenze von 200 Generationen und ein Zeitbudget von 600 Sekunden. Das dritte ist adaptiv und greift, wenn sich die beste Bewertung über 20 aufeinanderfolgende Generationen um weniger als ein Prozent verbessert hat. Die Kombination ist bewusst gewählt: Das adaptive Kriterium beendet den Lauf, sobald weitere Rechenzeit voraussichtlich keinen Gewinn mehr bringt, während die festen Grenzen sicherstellen, dass das Verfahren auch dann in vorgegebener Zeit ein Ergebnis liefert, wenn sich die Bewertung in kleinen Schritten immer weiter verbessert. Ein evolutionäres Verfahren liefert zu jedem Zeitpunkt eine gültige Lösung, sodass ein Abbruch nach Zeit stets ein verwertbares Ergebnis hinterlässt.

\section{Prototypische Umsetzung}

Kurzer Überblick über die prototypische Implementierung

\section{Fazit und Reflexion}
\label{sec:fazit}

% (1) Übernachtungen ohne Emissionen
%     In der Zielfunktion trägt nur der Transport zum CO2-Term bei. Die
%     Dekodierfunktion für die Hotelwahl (Abschnitt 2.3) kennt deshalb nur
%     w1 und w2, nicht w3. Real verursachen Übernachtungen sehr wohl
%     Emissionen, und zwar je nach Hotelkategorie unterschiedlich hohe.
%     Folge: Das Modell unterschätzt die Klimawirkung einer Reise
%     systematisch und kann keinen Zielkonflikt zwischen Hotelkomfort und
%     CO2 abbilden. Erweiterung wäre ein Emissionswert je Hotel und Nacht,
%     der den Decoder um den w3-Term ergänzt.
%
% (2) Feste Rundreise über alle Ziele
%     Pi ist eine Permutation, also wird jedes Ziel genau einmal besucht und
%     die Reise endet am Startort. Weglassen einzelner Ziele oder Mehrfach-
%     besuche sind nicht darstellbar, könnte aber durchaus relevant sein.
%
% (3) Zeitfenster statt konkreter Abfahrtszeiten
%     Die Diskretisierung in drei Fenster je Verkehrsmittel vereinfacht die
%     Nachbarschaftsstruktur, verwischt aber Unterschiede innerhalb eines
%     Fensters.
%
% (4) Statische Datenbasis
%     Alle Preise und Verbindungen stammen aus Nachschlagetabellen. Reale
%     Verfügbarkeiten und Preise ändern sich während der Planung.

% (5) Grenze der gewichteten Summe gegenueber Pareto
%     Die gewichtete Summe (Skalarisierung) findet je Gewichtung genau eine
%     Loesung, und zwar immer eine auf der konvexen Huelle der erreichbaren
%     Punktmenge. Pareto-optimale Plaene in einer Delle der Front sind fuer
%     KEIN Gewicht w optimal; beim Verschieben von w springt das Ergebnis
%     von einem Randpunkt zum anderen.
%     Gleiche Gewichte helfen NICHT. Gleiche Gewichte heissen "addiere beide Ziele gleich
%     stark", nicht "bevorzuge Ausgewogenheit". Eine Summe bewertet nur
%     den Gesamtwert, nicht die Verteilung auf die Teilziele. 
%
%     ANKNUEPFUNG: Die Abschlussbemerkung stellt die Frage, warum der
%     mittlere Bereich der Paretofront die hoechste Attraktivitaet besitzt.
%     Genau diesen Bereich verfehlt die gewichtete Summe systematisch.
%     QUELLE: Eiben/Smith S. 196 (Skalarisierung und ihre Nachteile) sowie
%     S. 197 (Nichtkonvexitaet und Diskontinuitaeten der Paretofront machen
%     Skalarisierungsverfahren Probleme). 

%     Ausblick: als Studienarbeit könnte man noch so und so ausbauen


% ===== LITERATURVERZEICHNIS =====
\newpage
\pagenumbering{Roman}
\setcounter{page}{\value{romanpages}}
\stepcounter{page}
\printbibliography[title={Literaturverzeichnis}, heading=bibintoc]

\end{document}
